\section{再利用する計算：特異部分への降下とBrauer周期}
\label{sec:descent-notes-v39}

偶数位数 $\ell=2d$ の中心chartについて、エタール被覆を降下する問題は
特異部分上で完全に解ける。$\Sigma$ の座標環は
$R=\mathbb C[Q^{\pm1},J]$、二単純加群の被覆は
\[
 L=R[\kappa]/(\kappa^2+tQ/\beta)
 \simeq\mathbb C[\kappa^{\pm1},J]
\]
である。$\iota(\kappa)=-\kappa$ とすると
\[
 A_\Sigma\simeq M_d\bigl(L\ltimes(L_\iota\oplus L_\iota)\bigr).
\]
平方零拡大の中心は $R$、根基を除いた商は $M_d(L)$ である。
これは閉点ごとの二頂点quiverに加え、二頂点が大域的には交換されることを記録する。

\subsection*{再利用できる座標と行列}

冪等元の降下には $x$ の固有値より $x^2$ の固有値を用いる。
\[
 z=-\frac{\beta x^2}{tQ},\qquad
 f_i=\frac1d\sum_{k=0}^{d-1}t^{-2ik}z^k,
 \qquad a_i=Q(t^{2i}-1).
\]
$a_i$ は $1\le i<d$ で単元であり、$f_i$ を結ぶ鎖から行列単位が得られる。
$f_0$ の角では $\kappa=xf_0$、$p=f_0u^df_0$、$r=f_0v^df_0$ で、
$p,r$ は $\kappa$ と反可換、全ての二次積が零である。
逆向きの表示に必要な係数は
\[
 c_d=\prod_{i=1}^{d-1}a_i=(-1)^{d-1}dQ^{d-1}.
\]
$x$ を $t^i\kappa$ の対角行列、$u$ を上向きの単位鎖と折返し $p$、
$v$ を逆向きの $a_i$ 倍の鎖と折返し $c_d^{-1}r$ とすれば、
元の代数の関係を満たし、$6d^2$ 個の自由基底を回収できる。

全ての単純加群を与える大域的表示族は $p=r=0$ で得られる。
被覆座標を内部的に読むには
\[
 \kappa=\frac{1-t}{2}\operatorname{Tr}(x)
\]
が有効である。$\operatorname{Tr}(x^d)$ は $d$ の偶奇で退化するので、
一様な座標としては使わない。
$Q$ の原点の周囲の一周は $\kappa\mapsto-\kappa$ であり、二単純加群を交換する。
$\lambda^{2d}=\Xi$ の被覆を $\Sigma$ に制限すると、
$Q=-t^{2i-1}\beta\lambda^2$ で定まる $d$ 個の連結成分に分かれ、各成分はこの二重被覆である。
特に $d$ が偶数のとき、単純加群の被覆は $\Xi$ の平方根ではない。

\subsection*{一般点の障害は全位数に有効}

一般点の中心体と代数は
\[
 K=\mathbb C(Q,J,\Xi,U),\qquad
 D=(\Xi,U)_t,\qquad
 x^\ell=\Xi,\quad u^\ell=U,\quad ux=txu
\]
である。整域性と $\dim_KD=\ell^2$ から指数は $\ell$。
$E=K(\Xi^{1/\ell})=\mathbb C(Q,J,\Xi^{1/\ell})(U)$ とし、
$U$ に関する付値を取れば
\[
 v_U(N_{E/K}(a))=\ell v_U(a).
\]
巡回代数のノルム判定により $[D]^r=0$ なら $\ell\mid r$、
また $U^\ell=N_{E/K}(U)$ なので周期も $\ell$ となる。
この付値の中心を現在のアフィンchart内の因子と同一視する必要はない。

同じBrauer類を持つ代数は、中心上の通常の自己準同型環
$\operatorname{End}_Z(M)$ とMorita同値にならない。
従って完成局所NCCRから大域NCCRへ移る際の障害が明示された。
有限エタール被覆 $\lambda^\ell=\Xi$ はAzumaya部分上でこの類を分裂させ、
位数 $\ell$ は一般点における最小有限分裂次数でもある。
分裂後に得られるのはベクトル束の自己準同型環であり、
その束が大域的に自由であるとは結論しない。

このBrauer計算は、$H$ に依存する別の中心chartを調べる際にも再利用できる。
一般点の記号代数が同じなら周期障害は同じであり、
$H$ の情報を区別するには中心の境界またはそのアフィンモデルを調べる必要がある。
